% from ega4-2-fr.tex
\oldpage[IV-2]{90}

\begin{proposition}[5.2.1]
\label{IV.2.5.2.1-fr}
Soient $k$ un corps, $X$ un préschéma irréductible localement de type fini sur
$k$, $\xi$ son point générique. Alors $X$ est biéquidimensionnel et l'on a
$\dim(X)=\operatorname{deg.tr}_k k(\xi)$.
\end{proposition}

\begin{proof}
Tout anneau local $\mathcal O_x$ est anneau local d'une $k$-algèbre de type
fini, donc quotient d'un anneau local régulier $(0,17.3.9)$, et l'on sait
$(0,16.5.12)$ que ces anneaux sont caténaires; par suite $X$ est un espace
caténaire (5.1.5) et comme $X$ est irréductible, il suffit (5.1.1) de montrer
que pour \emph{tout point fermé} $x\in X$, on a
\begin{equation}
\dim(\mathcal O_x)=\operatorname{deg.tr}_k k(\xi).
\tag{5.2.1.1}\label{IV.2.5.2.1.1-fr}
\end{equation}
On peut évidemment supposer $X$ réduit et affine, donc \emph{intègre}
d'anneau $A$, algèbre de type fini sur $k$. Soit
$n=\operatorname{deg.tr}_k k(\xi)$, $k(\xi)=K$ étant le corps des fractions
de $A$. On sait (Bourbaki, \emph{Alg. comm.}, chap.~V, \S~3,
n\textsuperscript{o}~1, th.~1) qu'il existe une sous-$k$-algèbre
$B=k[t_1,\ldots,t_n]$ de $A$, où les $t_i$ sont algébriquement indépendants
sur $k$, telle que $A$ soit une $B$-algèbre \emph{finie}. Soit
$\mathfrak m=\mathfrak j_x$, qui est par hypothèse un idéal maximal de $A$;
$\mathfrak n=B\cap\mathfrak m$ est donc un idéal maximal de $B$ (Bourbaki,
\emph{Alg. comm.}, chap.~V, \S~2, n\textsuperscript{o}~1, prop.~1), et
$A_{\mathfrak m}$ est un anneau local de la $B_{\mathfrak n}$-algèbre
\emph{finie} $S^{-1}A$, où $S=B-\mathfrak n$; comme $B_{\mathfrak n}$ est
intégralement clos et $S^{-1}A$ intègre, on a
$\dim(A_{\mathfrak m})=\dim(B_{\mathfrak n})$ $(0,16.1.6)$. On peut donc se
borner au cas où $A=k[t_1,\ldots,t_n]$; on sait alors que
$k'=k(x)$ est une extension \emph{finie} de $k$ (I, 6.4.2); soit
$A'=k'[t_1,\ldots,t_n]$; compte tenu de (I, 2.4.6 et 3.3.14), il existe un
idéal maximal $\mathfrak m'$ de $A'$ au-dessus de $\mathfrak m$ et tel que
$k'$ soit le corps résiduel de $A'_{\mathfrak m'}$. Comme $A'$ est intègre et
est une $A$-algèbre finie, le même raisonnement que ci-dessus montre que
$\dim(A'_{\mathfrak m'})=\dim(A_{\mathfrak m})$, si bien qu'on est ramené au
cas où $A/\mathfrak m=k$. Or dans ce cas $\mathfrak m$ est engendré par des
polynômes $t_i-a_i$ (où $a_i\in k$, $1\leq i\leq n$, les $a_i$ étant les
images canoniques des $t_i$ dans $A/\mathfrak m$). Remplaçant $t_i$ par
$t_i-a_i$, on voit finalement qu'on est ramené au cas où
$\mathfrak m=\sum_{i=1}^nAt_i$. Le complété de l'anneau local
$A_{\mathfrak m}$ est alors l'anneau de séries formelles
$k[[t_1,\ldots,t_n]]$; on sait qu'il a même dimension que $A_{\mathfrak m}$
$(0,16.2.4)$, et d'autre part la dimension de $k[[t_1,\ldots,t_n]]$ est égale
à $n$ $(0,17.1.4)$; d'où la conclusion.
\end{proof}

\begin{corollary}[5.2.2]
\label{IV.2.5.2.2-fr}
Pour un préschéma $X$ localement de type fini sur un corps $k$, $\dim(X)$
coïncide avec le nombre défini dans (4.1.1).
\end{corollary}

\begin{proof}
En effet, si $X_\alpha$ sont les sous-préschémas réduits ayant pour espaces
sous-jacents les composantes irréductibles de $X$, on a
$\dim(X)=\sup_\alpha(\dim(X_\alpha))$ $(0,14.1.2.1)$ et il suffit d'appliquer
(5.2.1) aux $X_\alpha$.
\end{proof}

\begin{corollary}[5.2.3]
\EGAReaderTarget{ega4.c04.cor.5-2-3}{IV.2.5.2.3-fr}
Soient $X$ un préschéma localement de type fini sur un corps $k$, $x$ un
point de $X$. On a
\begin{equation}
\dim_x(X)=\dim(\mathcal O_x)+\operatorname{deg.tr}_k k(x).
\tag{5.2.3.1}\label{IV.2.5.2.3.1-fr}
\end{equation}
\end{corollary}

\begin{proof}
On sait qu'il existe un voisinage ouvert $U$ de $x$ dans $X$ tel que
$\dim_x(X)=\dim(U)$ $(0,14.1.4.1)$, et l'on peut supposer que les composantes
irréductibles de $U$ soient les $X_i\cap U$, où les $X_i$ sont les
composantes irréductibles de $X$ contenant $x$; comme $U\cap X_i$

